% Options for packages loaded elsewhere
\PassOptionsToPackage{unicode}{hyperref}
\PassOptionsToPackage{hyphens}{url}
\PassOptionsToPackage{dvipsnames,svgnames,x11names}{xcolor}
%
\documentclass[
  10pt,
]{article}
\usepackage{amsmath,amssymb}
\usepackage{iftex}
\ifPDFTeX
  \usepackage[T1]{fontenc}
  \usepackage[utf8]{inputenc}
  \usepackage{textcomp} % provide euro and other symbols
\else % if luatex or xetex
  \usepackage{unicode-math} % this also loads fontspec
  \defaultfontfeatures{Scale=MatchLowercase}
  \defaultfontfeatures[\rmfamily]{Ligatures=TeX,Scale=1}
\fi
\usepackage{lmodern}
\ifPDFTeX\else
  % xetex/luatex font selection
\fi
% Use upquote if available, for straight quotes in verbatim environments
\IfFileExists{upquote.sty}{\usepackage{upquote}}{}
\IfFileExists{microtype.sty}{% use microtype if available
  \usepackage[]{microtype}
  \UseMicrotypeSet[protrusion]{basicmath} % disable protrusion for tt fonts
}{}
\makeatletter
\@ifundefined{KOMAClassName}{% if non-KOMA class
  \IfFileExists{parskip.sty}{%
    \usepackage{parskip}
  }{% else
    \setlength{\parindent}{0pt}
    \setlength{\parskip}{6pt plus 2pt minus 1pt}}
}{% if KOMA class
  \KOMAoptions{parskip=half}}
\makeatother
\usepackage{xcolor}
\usepackage[margin=0.9in]{geometry}
\usepackage{color}
\usepackage{fancyvrb}
\newcommand{\VerbBar}{|}
\newcommand{\VERB}{\Verb[commandchars=\\\{\}]}
\DefineVerbatimEnvironment{Highlighting}{Verbatim}{commandchars=\\\{\}}
% Add ',fontsize=\small' for more characters per line
\newenvironment{Shaded}{}{}
\newcommand{\AlertTok}[1]{\textcolor[rgb]{1.00,0.00,0.00}{\textbf{#1}}}
\newcommand{\AnnotationTok}[1]{\textcolor[rgb]{0.38,0.63,0.69}{\textbf{\textit{#1}}}}
\newcommand{\AttributeTok}[1]{\textcolor[rgb]{0.49,0.56,0.16}{#1}}
\newcommand{\BaseNTok}[1]{\textcolor[rgb]{0.25,0.63,0.44}{#1}}
\newcommand{\BuiltInTok}[1]{\textcolor[rgb]{0.00,0.50,0.00}{#1}}
\newcommand{\CharTok}[1]{\textcolor[rgb]{0.25,0.44,0.63}{#1}}
\newcommand{\CommentTok}[1]{\textcolor[rgb]{0.38,0.63,0.69}{\textit{#1}}}
\newcommand{\CommentVarTok}[1]{\textcolor[rgb]{0.38,0.63,0.69}{\textbf{\textit{#1}}}}
\newcommand{\ConstantTok}[1]{\textcolor[rgb]{0.53,0.00,0.00}{#1}}
\newcommand{\ControlFlowTok}[1]{\textcolor[rgb]{0.00,0.44,0.13}{\textbf{#1}}}
\newcommand{\DataTypeTok}[1]{\textcolor[rgb]{0.56,0.13,0.00}{#1}}
\newcommand{\DecValTok}[1]{\textcolor[rgb]{0.25,0.63,0.44}{#1}}
\newcommand{\DocumentationTok}[1]{\textcolor[rgb]{0.73,0.13,0.13}{\textit{#1}}}
\newcommand{\ErrorTok}[1]{\textcolor[rgb]{1.00,0.00,0.00}{\textbf{#1}}}
\newcommand{\ExtensionTok}[1]{#1}
\newcommand{\FloatTok}[1]{\textcolor[rgb]{0.25,0.63,0.44}{#1}}
\newcommand{\FunctionTok}[1]{\textcolor[rgb]{0.02,0.16,0.49}{#1}}
\newcommand{\ImportTok}[1]{\textcolor[rgb]{0.00,0.50,0.00}{\textbf{#1}}}
\newcommand{\InformationTok}[1]{\textcolor[rgb]{0.38,0.63,0.69}{\textbf{\textit{#1}}}}
\newcommand{\KeywordTok}[1]{\textcolor[rgb]{0.00,0.44,0.13}{\textbf{#1}}}
\newcommand{\NormalTok}[1]{#1}
\newcommand{\OperatorTok}[1]{\textcolor[rgb]{0.40,0.40,0.40}{#1}}
\newcommand{\OtherTok}[1]{\textcolor[rgb]{0.00,0.44,0.13}{#1}}
\newcommand{\PreprocessorTok}[1]{\textcolor[rgb]{0.74,0.48,0.00}{#1}}
\newcommand{\RegionMarkerTok}[1]{#1}
\newcommand{\SpecialCharTok}[1]{\textcolor[rgb]{0.25,0.44,0.63}{#1}}
\newcommand{\SpecialStringTok}[1]{\textcolor[rgb]{0.73,0.40,0.53}{#1}}
\newcommand{\StringTok}[1]{\textcolor[rgb]{0.25,0.44,0.63}{#1}}
\newcommand{\VariableTok}[1]{\textcolor[rgb]{0.10,0.09,0.49}{#1}}
\newcommand{\VerbatimStringTok}[1]{\textcolor[rgb]{0.25,0.44,0.63}{#1}}
\newcommand{\WarningTok}[1]{\textcolor[rgb]{0.38,0.63,0.69}{\textbf{\textit{#1}}}}
\setlength{\emergencystretch}{3em} % prevent overfull lines
\providecommand{\tightlist}{%
  \setlength{\itemsep}{0pt}\setlength{\parskip}{0pt}}
\setcounter{secnumdepth}{-\maxdimen} % remove section numbering
\usepackage{mathtools}
\usepackage{mathrsfs}
\setlength{\emergencystretch}{3em}
\ifLuaTeX
  \usepackage{selnolig}  % disable illegal ligatures
\fi
\usepackage{bookmark}
\IfFileExists{xurl.sty}{\usepackage{xurl}}{} % add URL line breaks if available
\urlstyle{same}
\hypersetup{
  colorlinks=true,
  linkcolor={Maroon},
  filecolor={Maroon},
  citecolor={Blue},
  urlcolor={Blue},
  pdfcreator={LaTeX via pandoc}}

\author{}
\date{}

\begin{document}

\section{ReLU population foundations for two-concept
SIM}\label{relu-population-foundations-for-two-concept-sim}

\textbf{Status:} Revision 2. The original foundations (P1--P9) are
retained, and Sections 15--24 add exact gradient attribution, the
matched linear control, immediate ReLU minor-entry generation, Gaussian
tail/purity bounds, a signed residual-curvature identity, and
gate-density/continuation interfaces. This remains a mathematical
research draft, not a proof of initialized competition (Theorem A), an
interval-arithmetic certificate, or a machine-checked formal proof.

\textbf{Scope:} Planar two-concept SIM; positive Gaussian noise; untied,
bias-free two-layer ReLU networks; both layers trained; aligned-balanced
isotropic initialization; population training data; fixed initial total
squared mass in the continuum limit.

\textbf{Preservation:} This is a new revision; the original foundations
files remain unchanged. It also leaves \texttt{SIM.md},
\texttt{main-15.tex}, and \texttt{ood\_reversals-appendix.tex}
untouched. The original SIM paper is a source for the task and
comparison model. The ReLU results below are derived here, not
attributed to that paper.

\subsection{Source map}\label{source-map}

\begin{itemize}
\tightlist
\item
  \textbf{{[}S1{]}} Yang, Park, Lubana, Okawa, Hu, and Tanaka,
  \emph{Swing-by Dynamics in Concept Learning and Compositional
  Generalization}, ICLR 2025, arXiv:2410.08309v2. The supplied PDF has
  32 pages. Relevant locations: Section 2.1, pp.~3--4; Section 4.2,
  pp.~7--9; Appendix D.1--D.2, pp.~19--25; Appendix E.4, p.~27.
\item
  \textbf{{[}S2{]}} \texttt{TodManlaibaatar/Swing-by-TwoLayerReLU},
  \texttt{full\_flow\_diagnostic.py}, functions \texttt{make\_data},
  \texttt{initialize}, \texttt{forward}, and \texttt{gradients}. The
  reviewed implementation uses a shared diagonal covariance, uniform
  planar input angles, \texttt{W\ =\ U.copy()}, and a loss equal to one
  half of the mean squared Euclidean residual norm.
\item
  \textbf{{[}S3{]}} Du, Hu, and Lee, \emph{Algorithmic Regularization in
  Learning Deep Homogeneous Models: Layers are Automatically Balanced},
  arXiv:1806.00900. Background for homogeneous-network balance
  identities; the particular identity needed here is proved directly
  below.
\end{itemize}

Reviewed code blob identifier:

\begin{Shaded}
\begin{Highlighting}[]
\NormalTok{0fccbe945ba9389609aa8bdcd8abdac1901d6ddd}
\end{Highlighting}
\end{Shaded}

No external convergence or mean-field theorem is used as a black box in
the proofs below. In particular, the global well-posedness statement is
for the labeled characteristic class specified here, not an unqualified
uniqueness claim for all distributional PDE solutions.

\section{1. What is inherited, and what must be
re-derived}\label{what-is-inherited-and-what-must-be-re-derived}

From {[}S1, Section 2.1{]}, inherit identity targets, equal cluster
weights, cluster means on distinct concept axes, and a common
coordinate-indexed covariance. Specializing to equal noise in two
dimensions is a restriction of that data model. We omit the optional
origin cluster.

Do not inherit the symmetric linear model's effective matrix evolution,
its tied-parameter metric, its discrete update, or its Appendix D
entrywise initialization assumptions. The current network has separate
trainable input and output weights. Its training field depends on actual
gates.

The original paper uses \texttt{s} for the number of concepts. Here
\texttt{s\ =\ d\ =\ 2}, and \texttt{S\_0} denotes initialization mass.
These symbols must not be interchanged.

The original compositional probe is the sum of concept means. Off-cone
probes are an extension of the evaluation family. We retain both
definitions without claiming that their error trajectories agree.

A useful proof strategy from {[}S1, Appendix D.2{]} is to establish
estimates up to a first exit and then close the corresponding bootstrap.
The controlled state and the signed estimates must nevertheless be
derived for ReLU.

\section{2. Locked setup and
normalization}\label{locked-setup-and-normalization}

\subsection{2.1 Data, architecture, and
loss}\label{data-architecture-and-loss}

In physical coordinates,

\[
 P_1^{\rm phys}=\mathcal N(\mu_1e_1,\sigma^2I_2),\qquad
 P_2^{\rm phys}=\mathcal N(\mu_2e_2,\sigma^2I_2),\qquad
 0<\mu_2<\mu_1,\quad\sigma>0.
\]

Let \(P^{\rm phys}=(P_1^{\rm phys}+P_2^{\rm phys})/2\). The finite-width
network and population-data loss are

\[
 f_h(x)=\sum_{i=1}^h w_i(u_i^\top x)_+,
 \qquad
 \mathcal L_h=\frac12\mathbb E_{P^{\rm phys}}\|f_h(X)-X\|^2.
 \tag{2.1}
\]

Both \(u_i\) and \(w_i\) are trained with the same unit-rate Euclidean
gradient-flow convention. No bias, weight decay, tied-weight constraint,
or frozen layer is present.

Use \textbf{target-minus-output} residuals throughout:

\[
 e_f(x)=x-f(x),\qquad E_\xi(f)=\frac12\|e_f(\xi)\|^2.
 \tag{2.2}
\]

The finite initialization is

\[
 \alpha_i\overset{\rm iid}{\sim}{\rm Unif}[0,2\pi),\qquad
 u_i(0)=w_i(0)=\sqrt{S_0/h}\,a_{\alpha_i},\qquad
 a_\theta=(\cos\theta,\sin\theta).
 \tag{2.3}
\]

Thus \(S_0=h\varepsilon_h^2\). The reviewed canonical implementation
{[}S2{]} has \(h=200\), \(\varepsilon_h=10^{-3}\), and therefore
\(S_0=2\cdot10^{-4}\). Keeping \(\varepsilon_h\) fixed as \(h\to\infty\)
is a different initialization regime and is not the limit defined here.

\subsection{2.2 Data scaling and time}\label{data-scaling-and-time}

Set

\[
 \rho=\mu_2/\mu_1,\qquad q=\sigma/\mu_1,\qquad
 \tau=\mu_1^2t_{\rm phys}.
 \tag{2.4}
\]

Positive homogeneity gives \(f(\mu_1x)=\mu_1 f(x)\). Consequently,

\[
 \mathcal L^{\rm phys}(\theta)=\mu_1^2\mathcal L^{\rm norm}(\theta),
 \qquad \theta^{\rm phys}(t)=\theta^{\rm norm}(\mu_1^2t).
 \tag{2.5}
\]

There is no accompanying rescaling of \(S_0\). From now on, \(t\)
denotes normalized time and

\[
 P_1=\mathcal N(e_1,q^2I_2),\quad
 P_2=\mathcal N(\rho e_2,q^2I_2),\quad
 P=(P_1+P_2)/2.
 \tag{2.6}
\]

Write \(c_1=e_1\), \(c_2=\rho e_2\). Define

\[
 M_{2,p}=\mathbb E_{P_p}|X|^2=|c_p|^2+2q^2,\qquad
 M_2=\mathbb E_P|X|^2=\frac{1+\rho^2}{2}+2q^2.
 \tag{2.7}
\]

The uncentered second-moment matrix is

\[
 A=\mathbb E_P[XX^\top]
 =\operatorname{diag}(1/2+q^2,\rho^2/2+q^2).
 \tag{2.8}
\]

For a fixed numerical probe \(\xi\), physical-time rates are \(\mu_1^2\)
times the corresponding normalized-time rates. For the additionally
rescaled probe \(\xi^{\rm phys}=\mu_1\bar\xi\), errors scale by
\(\mu_1^2\) and physical-time rates by \(\mu_1^4\). Never reuse a
physical time from the canonical experiment without this conversion.

\subsection{2.3 Exact link to averaged particle
coordinates}\label{exact-link-to-averaged-particle-coordinates}

Set \(U_i=\sqrt h\,u_i\), \(W_i=\sqrt h\,w_i\). Then

\[
 f_h(x)=\frac1h\sum_iW_i(U_i^\top x)_+,
\]

but their equations in the same physical/normalized time are

\[
 \dot W_i=\mathbb E_P[e_f(X)(U_i^\top X)_+],\qquad
 \dot U_i=\mathbb E_P[X\,e_f(X)^\top W_i\,\mathbf1_{\{U_i^\top X>0\}}].
 \tag{2.9}
\]

These follow by multiplying the original equations by \(\sqrt h\).
Equivalently, the averaged network uses the mean-field metric, or
learning rate \(h\) relative to ordinary Euclidean gradient descent in
\((U_i,W_i)\). Equation (2.9), not an unspecified convention, defines
the continuum time normalization.

\section{3. Proposition P1: exact gradient field and
balance}\label{proposition-p1-exact-gradient-field-and-balance}

For a network output \(f\) and an input direction \(a_\theta\), define
the cluster-conditioned residual matrix

\[
 R_{p,f}(\theta)
 =\mathbb E_{P_p}[e_f(X)X^\top\mathbf1_{\{a_\theta^\top X>0\}}],
 \qquad R_f=(R_{1,f}+R_{2,f})/2.
 \tag{3.1}
\]

For every nonzero input weight \(u_i\), population gradient flow is

\[
 \dot w_i=R_f(\arg u_i)u_i,\qquad
 \dot u_i=R_f(\arg u_i)^\top w_i.
 \tag{3.2}
\]

Moreover,

\[
 \frac d{dt}(|w_i|^2-|u_i|^2)=0.
 \tag{3.3}
\]

\textbf{Proof.} For fixed nonzero \(u_i\), the Gaussian law assigns zero
probability to \(u_i^\top X=0\). On bounded parameter sets, the
derivatives of the integrand are bounded by a constant times \(|X|^2\),
which is integrable. Differentiating the loss under the expectation
therefore gives

\[
 -\nabla_{w_i}\mathcal L=\mathbb E[e_f(X)(u_i^\top X)_+],
\]

\[
 -\nabla_{u_i}\mathcal L=\mathbb E[X(w_i^\top e_f(X))\mathbf1_{\{u_i^\top X>0\}}].
\]

Use \((u_i^\top X)_+=(u_i^\top X)\mathbf1_{\{u_i^\top X>0\}}\) to obtain
(3.2). Then

\[
 \frac d{dt}|w_i|^2=2w_i^\top R_fu_i
 =2u_i^\top R_f^\top w_i=\frac d{dt}|u_i|^2.
\]

This proves (3.3). Nonvanishing on every finite horizon for the
specified initialized flow follows from P5 below. \(\square\)

\textbf{Balance does not imply persistent alignment.} At a state with
\(u_i=w_i\),

\[
 \dot w_i-\dot u_i=(R_f-R_f^\top)u_i.
\]

The right side is not identically zero for a general state. We keep two
angles even though they start equal.

\section{4. Proposition P2: exact angular state and transport--reaction
flow}\label{proposition-p2-exact-angular-state-and-transportreaction-flow}

Let \(\lambda(d\alpha)=d\alpha/(2\pi)\) on the initial-angle circle. A
balanced labeled continuum neuron is written

\[
 U_t(\alpha)=\sqrt{m_t(\alpha)}a_{\theta_t(\alpha)},\qquad
 W_t(\alpha)=\sqrt{m_t(\alpha)}b_{\psi_t(\alpha)},
\]

where \(b_\psi=(\cos\psi,\sin\psi)\) is an output direction and
\(m_t>0\). Set

\[
 \nu_t=\int m_t(\alpha)
 \delta_{(\theta_t(\alpha),\psi_t(\alpha))}\,d\lambda(\alpha),
 \qquad S(t)=\nu_t(\mathbb T^2)=\int m_t\,d\lambda.
 \tag{4.1}
\]

Then exactly

\[
 f_{\nu_t}(x)=\int b_\psi(a_\theta^\top x)_+\,d\nu_t(\theta,\psi).
 \tag{4.2}
\]

For \(R=R_{f_{\nu_t}}(\theta)\), define

\[
 g_\nu(\theta,\psi)=b_\psi^\top R a_\theta,\qquad
 v_\theta=b_\psi^\top R a_\theta^\perp,\qquad
 v_\psi=(b_\psi^\perp)^\top R a_\theta,
 \tag{4.3}
\]

where \(a_\theta^\perp=(-\sin\theta,\cos\theta)\) and similarly for
\(b_\psi^\perp\). The exact characteristic equations are

\[
 \boxed{\dot\theta=v_\theta,\qquad
 \dot\psi=v_\psi,\qquad \dot m=2mg_\nu.}
 \tag{4.4}
\]

The initial condition is

\[
 \theta_0(\alpha)=\psi_0(\alpha)=\alpha,\qquad m_0(\alpha)=S_0,
\]

\[
 \nu_0=\frac{S_0}{2\pi}\int_0^{2\pi}\delta_{(\alpha,\alpha)}\,d\alpha.
 \tag{4.5}
\]

For every periodic \(C^1\) test function \(\varphi(\theta,\psi)\),

\[
 \frac d{dt}\int\varphi\,d\nu_t
 =\int(v_\theta\partial_\theta\varphi+
 v_\psi\partial_\psi\varphi+2g_\nu\varphi)\,d\nu_t.
 \tag{4.6}
\]

Thus, in weak form,

\[
 \partial_t\nu+\partial_\theta(v_\theta\nu)
 +\partial_\psi(v_\psi\nu)=2g_\nu\nu.
 \tag{4.7}
\]

\textbf{Proof.} Substitute \(U=\sqrt m\,a\), \(W=\sqrt m\,b\) into
(2.9). The projections of \(\dot U=\sqrt m R^\top b\) onto \(a\) and
\(a^\perp\) yield \(\dot m/(2m)=b^\top Ra\) and
\(\dot\theta=b^\top Ra^\perp\). Projections of \(\dot W=\sqrt m Ra\)
onto \(b\) and \(b^\perp\) give the same radial equation and the
displayed output-angle equation. This proves (4.4). The network
representation follows from positive homogeneity. Differentiating
\(m_t(\alpha)\varphi(\theta_t(\alpha),\psi_t(\alpha))\) and integrating
proves (4.6). All differentiations are justified locally by bounded
characteristic variables; P5 gives the bounds on every finite horizon.
\(\square\)

The measure need not have a density on \(\mathbb T^2\). Here it starts
on, and is represented by the evolution of, a one-dimensional labeled
curve. Angular spread, input/output mismatch, and all label masses are
retained. Projection of this curve onto the input-angle coordinate need
not stay one-to-one.

For a finite network, the identical angular representation uses

\[
 \nu_h=\sum_{i=1}^h |u_i|^2\delta_{(\arg u_i,\arg w_i)}.
\]

The continuum equations are therefore not a two-block closure. A
quantitative random finite-width/sample approximation theorem is a
separate task.

\section{5. Proposition P3: exact initial output and learning
values}\label{proposition-p3-exact-initial-output-and-learning-values}

For the initial measure (4.5),

\[
 f_0(x)=\frac{S_0}{4}x,\qquad
 \int b_\psi a_\theta^\top\,d\nu_0=\frac{S_0}{2}I_2.
 \tag{5.1}
\]

Consequently, with \(c_0=1-S_0/4\),

\[
 e_0(x)=c_0x,\qquad
 E_\xi(0)=\frac{c_0^2}{2}|\xi|^2,
\]

\[
 \mathcal L(0)=\frac{c_0^2}{2}M_2,\qquad
 \mathcal L_p(0)=\frac{c_0^2}{2}M_{2,p}.
 \tag{5.2}
\]

\textbf{Proof.} Pair \(a\) and \(-a\) under the uniform angular law:

\[
 a(a^\top x)_++(-a)((-a)^\top x)_+=aa^\top x.
\]

It follows that
\(\mathbb E[a(a^\top x)_+]=\frac12\mathbb E[aa^\top]x=x/4\), since the
diagonal angular averages are \(1/2\) and the off-diagonal averages
vanish. Multiplication by \(S_0\) gives the first identity. The same
angular second moment gives the second. Substitute the resulting
residual into the losses. \(\square\)

The finite random network does not equal its angular average exactly.
Equation (5.1) is an exact continuum initialization identity, not an
exact statement about every finite random draw. In particular, setting
\(f_0=0\) would change the initial problem. Also, the effective ungated
moment in (5.1) has zero off-diagonal entries; {[}S1, Eq. (D.8){]} must
not be imported as an initialization assumption for it.

\section{6. Proposition P4: exact one-dimensional Gaussian
calculus}\label{proposition-p4-exact-one-dimensional-gaussian-calculus}

\subsection{6.1 Radial integration}\label{radial-integration}

Let \(p_p(x)\) be the planar density of \(P_p\) and define a weighted
angular density, with respect to unnormalized Lebesgue measure
\(d\beta\),

\[
 \omega_p(\beta)=\int_0^\infty r^3p_p(r a_\beta)\,dr,
 \qquad \omega=(\omega_1+\omega_2)/2.
 \tag{6.1}
\]

For every integrable positively 2-homogeneous scalar, vector, or matrix
function \(H\),

\[
 \mathbb E_{P_p}H(X)=\int_0^{2\pi}\omega_p(\beta)H(a_\beta)\,d\beta.
 \tag{6.2}
\]

In particular \(\int\omega_p=M_{2,p}\) and \(\int\omega=M_2\). These
angular weights are not probability densities; do not insert an
additional \(1/(2\pi)\).

\textbf{Proof.} Write \(X=r a_\beta\) in the planar integral, whose
Jacobian is \(r\). Since \(H(r a_\beta)=r^2H(a_\beta)\) for \(r\ge0\),
its radial weight is precisely \(r^3p_p(r a_\beta)\). Fubini yields
(6.2). \(\square\)

The residual is positively 1-homogeneous because both the target and
network are. Thus every training integral needed here has degree two,
including loss, residual matrices, feature kernels, and the Gaussian
expectation in the probe-rate formula. No Gaussian-tail truncation has
been made.

\subsection{6.2 Explicit angular weight}\label{explicit-angular-weight}

Let \(\ell=c_p^\top a_\beta\). Then

\[
 \omega_p(\beta)=
 \frac{\exp[-(|c_p|^2-\ell^2)/(2q^2)]}{2\pi q^2}J_3(\ell,q),
 \quad
 J_3(\ell,q)=\int_0^\infty r^3e^{-(r-\ell)^2/(2q^2)}\,dr.
 \tag{6.3}
\]

Writing \(J_0=q\sqrt{2\pi}\,\Phi(\ell/q)\), integration by parts gives

\[
 J_1=\ell J_0+q^2e^{-\ell^2/(2q^2)},\qquad
 J_k=\ell J_{k-1}+(k-1)q^2J_{k-2}\quad(k\ge2),
\]

and hence

\[
 J_3=(\ell^3+3\ell q^2)J_0+
 q^2(\ell^2+2q^2)e^{-\ell^2/(2q^2)}.
 \tag{6.4}
\]

The integral (6.3) is nonnegative and remains an appropriate definition
when (6.4) has numerical cancellation for negative \(\ell/q\).

For \(q>0\), \(\omega_p\) is smooth and bounded. An explicit exact
supremum is

\[
 W_p:=\|\omega_p\|_\infty=
 \frac{J_3(|c_p|,q)}{2\pi q^2}.
 \tag{6.5}
\]

Indeed the density \(p_p(r a_\beta)\) is, for every \(r\ge0\), maximized
when \(a_\beta\) points along \(c_p\). Put
\(W=\|\omega\|_\infty\le(W_1+W_2)/2\).

\subsection{6.3 Gaussian feature
kernels}\label{gaussian-feature-kernels}

Define

\[
 k_\theta(\beta)=(a_\theta^\top a_\beta)_+,
\]

\[
 B_p(\theta)=\mathbb E_{P_p}[X(a_\theta^\top X)_+]
 =\int\omega_p(\beta)a_\beta k_\theta(\beta)\,d\beta,
\]

\[
 C_p(\theta,\vartheta)
 =\mathbb E_{P_p}[(a_\theta^\top X)_+(a_\vartheta^\top X)_+]
 =\int\omega_p k_\theta k_\vartheta\,d\beta.
 \tag{6.6}
\]

Let \(\ell=a_\theta^\top c_p\), \(\kappa=\ell/q\), and let \(\phi,\Phi\)
be the standard normal density and CDF. Then

\[
 B_p(\theta)=c_p[\ell\Phi(\kappa)+q\phi(\kappa)]
 +q^2a_\theta\Phi(\kappa).
 \tag{6.7}
\]

\textbf{Proof of (6.7).} Decompose
\(X=c_p+qZa_\theta+qZ_\perp a_\theta^\perp\) with independent standard
normal variables. The perpendicular centered term has mean zero. Direct
one-dimensional integration gives
\(\mathbb E(\ell+qZ)_+=\ell\Phi(\kappa)+q\phi(\kappa)\) and
\(\mathbb E[Z(\ell+qZ)_+]=q\Phi(\kappa)\). Substitution proves the
formula. \(\square\)

The residual vector field admits the exact kernel closure

\[
 G_{p,\nu}(\theta):=R_{p,\nu}(\theta)a_\theta
 =B_p(\theta)-\int b_\varphi C_p(\theta,\vartheta)\,d\nu(\vartheta,\varphi).
 \tag{6.8}
\]

Differentiation is with respect to the displayed test angle, holding
\(\nu\) fixed:

\[
 \partial_\theta G_{p,\nu}(\theta)
 =R_{p,\nu}(\theta)a_\theta^\perp.
 \tag{6.9}
\]

There is no boundary term here because the feature \(k_\theta\) itself
vanishes at a gate boundary.

First derivatives of \(B_p\) and \(C_p\) exist by differentiation under
the angular integral. Useful explicit second-derivative formulas are

\[
 \partial_\theta^2B_p(\theta)=-B_p(\theta)
 +\sum_{\epsilon\in\{-1,1\}}
 \omega_p(\theta+\epsilon\pi/2)a_{\theta+\epsilon\pi/2},
\]

\[
 \partial_\theta^2C_p(\theta,\vartheta)=-C_p(\theta,\vartheta)
 +\sum_{\epsilon\in\{-1,1\}}
 \omega_p(\theta+\epsilon\pi/2)k_\vartheta(\theta+\epsilon\pi/2).
 \tag{6.10}
\]

The mixed derivative is the integral of
\(\omega_p(\partial_\theta k_\theta)(\partial_\vartheta k_\vartheta)\).
All these derivatives are continuous: the boundary terms are continuous,
and dominated convergence handles the mixed derivative. In particular,
\(C_p\) is jointly \(C^2\) and its first derivatives are Lipschitz on
the angular torus. Its pure second derivatives have magnitude at most
\(M_{2,p}+2W_p\), and its mixed derivative at most \(M_{2,p}\). No
higher-order analyticity is asserted.

\subsection{6.4 A simpler regularity formula for the residual
matrix}\label{a-simpler-regularity-formula-for-the-residual-matrix}

Since \(|f_\nu(x)|\le S|x|\), the function \(e_\nu(a_\beta)\) is
continuous and bounded by \(1+S\). Equation (6.2) gives

\[
 R_{p,\nu}(\theta)=
 \int_{\theta-\pi/2}^{\theta+\pi/2}
 \omega_p(\beta)e_\nu(a_\beta)a_\beta^\top\,d\beta.
 \tag{6.11}
\]

Use periodic integration in this formula. Therefore \(R_{p,\nu}\) is
\(C^1\) in \(\theta\), with

\[
 \partial_\theta R_{p,\nu}(\theta)=
 [\omega_p e_\nu(a_\beta)a_\beta^\top]_{\beta=\theta+\pi/2}
 -[\omega_p e_\nu(a_\beta)a_\beta^\top]_{\beta=\theta-\pi/2}.
 \tag{6.12}
\]

Consequently,

\[
 \|R_{p,\nu}(\theta)\|\le M_{2,p}(1+S),\qquad
 \|\partial_\theta R_{p,\nu}(\theta)\|\le2W_p(1+S).
 \tag{6.13}
\]

All matrix norms are operator norms. The same bounds hold for the
mixture using \(M_2,W\). This proves the training-field regularity
directly; it does not assume a density for the neuron input-angle
marginal.

\section{7. Proposition P5: global characteristic existence,
dissipation, and
bounds}\label{proposition-p5-global-characteristic-existence-dissipation-and-bounds}

For every fixed \(0<\rho<1\), \(q>0\), and \(S_0>0\), the initial
condition (4.5) has a unique global solution in the continuous
labeled-characteristic class. Its masses remain strictly positive on
every finite horizon.

The population loss satisfies

\[
 \dot{\mathcal L}(t)
 =-\int\bigl(\|R_\nu a_\theta\|^2+
                  \|R_\nu^\top b_\psi\|^2\bigr)\,d\nu_t
 =-\int(2g_\nu^2+v_\theta^2+v_\psi^2)\,d\nu_t\le0.
 \tag{7.1}
\]

Total mass satisfies the exact identity

\[
 \boxed{
 \dot S(t)=2\mathbb E_P[e_t(X)^\top f_t(X)]
 =\frac{M_2}{2}-2\mathbb E_P\left|f_t(X)-\frac X2\right|^2
 \le\frac{M_2}{2}.}
 \tag{7.2}
\]

Set

\[
 C_0:=\sqrt{2\mathcal L(0)M_2}=|1-S_0/4|M_2.
 \tag{7.3}
\]

On every finite horizon,

\[
 |g_\nu|,\ |v_\theta|,\ |v_\psi|\le C_0,
 \qquad
 S_0e^{-2C_0t}\le m_t(\alpha)\le S_0e^{2C_0t},
 \tag{7.4}
\]

\[
 0<S(t)\le\min\{S_0e^{2C_0t},\ S_0+M_2t/2\}.
 \tag{7.5}
\]

Angles in continuous lifts move by at most \(C_0t\) each. Further,

\[
 |f_t(x)|\le S(t)|x|,\qquad
 |\dot f_t(x)|\le2C_0S(t)|x|\quad\text{for a.e. }t
 \tag{7.6}
\]

for each fixed \(x\).

\textbf{Proof: local existence.} Use the variables
\((\theta(\alpha)-\alpha,\psi(\alpha)-\alpha,m(\alpha))\) in the Banach
space of continuous periodic functions with the sup norm, restricting
initially to masses bounded away from zero. For two labeled states with
mass and sup-norm bounds, positive homogeneity and the Lipschitz
property of ReLU give a uniform bound on the difference of their outputs
on the unit circle. Integrating this difference in (6.11) bounds the
change in \(R\) at the same angle. Equation (6.13) bounds its change
when the angle moves. Trigonometric factors are Lipschitz, and
multiplication by bounded \(m\) is Lipschitz. Hence the right side of
(4.4) is locally Lipschitz in this Banach space. The usual local ODE
construction applies: the integral map
\(Z\mapsto Z_0+\int_0^t\mathscr F(Z_s)ds\) is a contraction on a
sufficiently small closed tube because its Lipschitz constant is at most
the field constant times the interval length. Its fixed point is the
unique local solution. The detailed estimates also appear in P8.

\textbf{Proof: dissipation on a local interval.} Define
\(U=\sqrt m\,a\), \(W=\sqrt m\,b\). Their velocities are exactly
\(\dot W=RU\) and \(\dot U=R^\top W\). The derivative of
\(W(U^\top X)_+\) exists for almost every training \(X\) for each label.
Gaussian boundary-nullness and Fubini justify differentiation under the
training and label integrals, with a dominating multiple of \(|X|^2\) on
the local bounded tube. Thus

\[
 \dot{\mathcal L}
 =-\mathbb E[e_f(X)^\top\dot f(X)]
 =-\int (|\dot W|^2+|\dot U|^2)\,d\lambda.
\]

Substitute the polar forms and decompose \(Ra\) along \(b,b^\perp\) and
\(R^\top b\) along \(a,a^\perp\). Their shared radial component is
\(g\). This gives (7.1).

\textbf{Proof: mass and continuation.} Integrate \(\dot m=2mg\) and use
the definitions to get \(\dot S=2\mathbb E[e^\top f]\). Completing the
square gives (7.2). Also

\[
 \|R_\nu(\theta)\|
 \le\mathbb E[|e_f(X)||X|]
 \le\sqrt{2\mathcal L(t)M_2}\le C_0.
\]

Equations (4.4) imply (7.4), and integration of (7.2) gives (7.5). These
bounds keep the characteristic variables bounded and masses bounded away
from zero on any prescribed finite interval. Therefore the local
solution can be continued past any finite endpoint by the same local
construction. Uniqueness is preserved. Equation (7.6) follows by
differentiating a single feature and using
\(\|Ra\|,\|R^\top b\|\le C_0\); each of the two terms has magnitude at
most \(mC_0|x|\). \(\square\)

These are structural, potentially conservative bounds. In particular,
global existence and bounded mass do not imply that the trajectory
enters the competition region or that the bounds are sharp enough to
certify the desired late sign margins.

\section{8. Proposition P6: exact cluster-sourced probe
rates}\label{proposition-p6-exact-cluster-sourced-probe-rates}

For \(z=(\theta,\psi)\) write
\(I_\theta(x)=\mathbf1_{\{a_\theta^\top x>0\}}\), taking value zero at
equality. Define the matrix-valued tangent kernel

\[
 K_\nu(\xi,x)=\int\left[
 (a_\theta^\top\xi)_+(a_\theta^\top x)_+I_2
 +(\xi^\top x)I_\theta(\xi)I_\theta(x)b_\psi b_\psi^\top
 \right]d\nu(z).
 \tag{8.1}
\]

The raw cluster-\(p\) probe velocity is

\[
 V_{p,\nu}(\xi)=\mathbb E_{P_p}[K_\nu(\xi,X)e_\nu(X)].
 \tag{8.2}
\]

An equivalent formula, convenient for direct evaluation, is

\[
 V_{p,\nu}(\xi)=\int\left[
 R_{p,\nu}(\theta)a_\theta(a_\theta^\top\xi)_+
 +I_\theta(\xi)b_\psi b_\psi^\top R_{p,\nu}(\theta)\xi
 \right]d\nu.
 \tag{8.3}
\]

Define

\[
 \boxed{D_p(t,\xi)=-\frac12e_t(\xi)^\top V_{p,\nu_t}(\xi).}
 \tag{8.4}
\]

For each fixed probe \(\xi\),

\[
 \dot f_t(\xi)=\frac12[V_{1,\nu_t}(\xi)+V_{2,\nu_t}(\xi)],\qquad
 \boxed{\dot E_\xi(t)=D_1(t,\xi)+D_2(t,\xi)}
 \tag{8.5}
\]

for almost every time. Equivalently, writing the output-minus-target
residual as \(r=f-x\) only in the following identity,

\[
 D_p=-\frac12\mathbb E_{P_p}
 [\langle r_t(\xi),K_{\nu_t}(\xi,X)r_t(X)\rangle].
 \tag{8.6}
\]

\textbf{Proof.} Differentiate \(W(U^\top\xi)_+\) using (2.9). The
output-weight derivative contributes
\(m R_\nu a_\theta(a_\theta^\top\xi)_+\). The input-weight derivative
contributes \(mI_\theta(\xi)b_\psi b_\psi^\top R_\nu\xi\). Integrate
labels and substitute \(R_\nu=(R_1+R_2)/2\) to obtain (8.3)--(8.5).
Expanding \(R_p\) and interchanging integrals gives (8.1)--(8.2).
Finally \(\dot E_\xi=-e_\xi^\top\dot f_\xi\), proving (8.4)--(8.6).

For completeness, the a.e. statement at probe gates does not require
training-density arguments. For each label, \(z(t)=U_t(\alpha)^\top\xi\)
is absolutely continuous; \(d(z_+)/dt=\mathbf1_{\{z>0\}}\dot z\) for
a.e. \(t\), with \(\dot z=0\) for a.e. time in the zero set. Applying
Fubini in time and label and the finite-horizon bound from P5 justifies
the integrated identity for each fixed probe. \(\square\)

\textbf{Boundary convention.} When \(\nu\{a_\theta^\top\xi=0\}=0\), the
cluster-rate formulas are independent of the chosen value of the
indicator at zero. If positive mass sits on a probe boundary, the
displayed zero-gate convention still gives the a.e. total identity, but
individual cluster attributions at that boundary can depend on the
convention. A theorem asserting robust common-interval signs must
control this issue; see P9.

\textbf{Interpretation.} The index \(p\) specifies the training
distribution sourcing the gradient. Every neuron contributes to both
\(D_p\) according to its actual gates. No strong/weak neuron partition,
deletion, frozen gate, or retraining intervention is part of these
definitions.

The kernel is positive semidefinite as a kernel: for any
square-integrable vector test function \(h\) for which the expressions
exist,

\[
 \iint h(x)^\top K_\nu(x,y)h(y)\,dP(x)dP(y)
 =\int\left\{\left|\mathbb E[h(X)(a_\theta^\top X)_+]\right|^2
 +\left|\mathbb E[X(b_\psi^\top h(X))I_\theta(X)]\right|^2\right\}d\nu\ge0.
\]

This does not fix the sign of an individual off-diagonal probe/training
contraction \(D_p\).

\section{9. Proposition P7: exact concept-learning
observables}\label{proposition-p7-exact-concept-learning-observables}

Let \([B_p(\theta)]_p\) denote coordinate \(p\) of \(B_p\). Define

\[
 \mathsf m_p(\nu)=
 \frac{\mathbb E_{P_p}[X_p f_{\nu,p}(X)]}{\mathbb E_{P_p}[X_p^2]}
 =\frac{\int (b_\psi)_p[B_p(\theta)]_p\,d\nu}
 {|c_p|^2+q^2}.
 \tag{9.1}
\]

Then \(\mathsf m_p(\nu_0)=S_0/4\). The cluster loss is exactly

\[
 \mathcal L_p(\nu)=\frac{M_{2,p}}2
 -\int b_\psi^\top B_p(\theta)\,d\nu(\theta,\psi)
\]

\[
 \hspace{8mm}+\frac12\iint
 b_\psi^\top b_\varphi C_p(\theta,\vartheta)
 \,d\nu(\theta,\psi)d\nu(\vartheta,\varphi),
 \qquad \mathcal L=\frac{\mathcal L_1+\mathcal L_2}{2}.
 \tag{9.2}
\]

\textbf{Proof.} Substitute (4.2) into the correlation defining
\(\mathsf m_p\), then use (6.6). For the loss, expand \(|X-f(X)|^2\)
into \(|X|^2-2X^\top f(X)+|f(X)|^2\) and evaluate the two network terms
by Fubini and (6.6). At initialization use \(f_0=(S_0/4)\mathrm{id}\).
\(\square\)

A future substantial-learning conclusion should use both
\(\mathsf m_2\ge c_{\rm weak}>0\) and a nontrivial relative cluster-loss
reduction, rather than reuse a strong block's small transverse entry
under the ambiguous symbol \(m_{22}\).

Neither monotonicity of \(\mathsf m_1\) nor monotonicity of an
individual \(\mathcal L_p\) follows from P5. Only the averaged training
loss is known to decrease. Any progress-clock inequality is a later
trajectory estimate, not a foundational fact.

\section{10. Proposition P8: finite-horizon labeled-state and data
stability}\label{proposition-p8-finite-horizon-labeled-state-and-data-stability}

\subsection{10.1 State stability for the same data
law}\label{state-stability-for-the-same-data-law}

Compare two characteristic solutions on a common label space, with
corresponding continuous angular lifts and masses. Suppose
\(S,\widetilde S\le M\) on \([0,T]\). Define

\[
 A_\theta=\|\theta-\widetilde\theta\|_\infty,\qquad
 A_\psi=\|\psi-\widetilde\psi\|_\infty,\qquad
 A_*=A_\theta+A_\psi,
\]

\[
 B_*=\int|m-\widetilde m|\,d\lambda,\qquad
 \Delta=A_*+B_*,\qquad \eta_f=B_*+MA_*.
 \tag{10.1}
\]

Then

\[
 \sup_{|x|=1}|f_\nu(x)-f_{\widetilde\nu}(x)|\le\eta_f.
 \tag{10.2}
\]

Let

\[
 H=M_2(1+M),\qquad J=2W(1+M),\qquad
 F=H+M_2M+J,
\]

\[
 L_M=\max\{2H+2(1+M)M_2,\ 2(1+M)F\}.
 \tag{10.3}
\]

Then

\[
 \Delta(t)\le e^{L_Mt}\Delta(0),\qquad 0\le t\le T.
 \tag{10.4}
\]

\textbf{Proof.} Split the network difference into mass,
output-direction, and feature differences. Since \(|b|\le1\) and ReLU is
1-Lipschitz, each unit-input feature has norm at most one, and the
latter two differences cost at most \(\widetilde S(A_\psi+A_\theta)\).
This proves (10.2).

At corresponding test angles,

\[
 \|R_\nu(\theta)-R_{\widetilde\nu}(\widetilde\theta)\|
 \le M_2\eta_f+J A_\theta.
 \tag{10.5}
\]

The first term follows by integrating the residual difference against
\(|X|\); the second follows from (6.13). Therefore each of the
differences of \(g,v_\theta,v_\psi\) is at most

\[
 Q:=M_2 B_*+F A_*.
\]

The upper right derivative of \(A_*\) is at most \(2Q\). The mass
equation gives

\[
 \frac{d^+}{dt}B_*
 \le2H B_*+2\int\widetilde m|g-\widetilde g|\,d\lambda
 \le2H B_*+2M Q.
\]

Add the inequalities to obtain \(\Delta'\le L_M\Delta\) in the
upper-derivative/integral sense. Multiplication by \(e^{-L_Mt}\) and
integration proves (10.4). The same bounds with a local sup bound on
\(m\) establish the Banach-space local Lipschitz assertion used in P5.
\(\square\)

This is a labeled coupling bound, not a total-variation estimate between
singular angular measures. Its constants are explicit but need not be
sharp enough for canonical numerical certification.

\subsection{10.2 Perturbing Gaussian
parameters}\label{perturbing-gaussian-parameters}

Let the second law have angular weight \(\widetilde\omega\). Put

\[
 \delta_\omega=\|\omega-\widetilde\omega\|_{L^1(d\beta)}.
\]

Choose \(H,J,F,L_M\) using uniform upper bounds valid for both laws.
Equation (10.5) acquires the extra term \((1+M)\delta_\omega\), so

\[
 \Delta(t)\le e^{L_Mt}\Delta(0)
 +2(1+M)^2\delta_\omega\frac{e^{L_Mt}-1}{L_M}.
 \tag{10.6}
\]

On a compact parameter rectangle with \(q\ge q_{\min}>0\), let \(M_4^*\)
bound

\[
 \mathbb E_{\mathcal N(c,q^2I_2)}|X|^4
 =|c|^4+8q^2|c|^2+8q^4
\]

for both cluster means and every point of the rectangle. Then a valid
bound is

\[
 \delta_\omega\le
 \frac{\sqrt{M_4^*}}{q_{\min}}
 \bigl(|\rho-\widetilde\rho|+2|q-\widetilde q|\bigr).
 \tag{10.7}
\]

\textbf{Proof.} At a fixed network state, the difference between the
residual matrices for the two angular weights is bounded by
\((1+M)\delta_\omega\), using (6.11). Add this to the proof of (10.4)
and integrate the resulting scalar differential inequality to get
(10.6).

For (10.7), differentiate the planar Gaussian density. Its score for a
unit mean displacement is \((v^\top Z)/q\); its score for \(q\) is
\((|Z|^2-2)/q\), with \(Z\sim\mathcal N(0,I_2)\). Integrating absolute
scores with weight \(|X|^2\) bounds the corresponding \(L^1\)
derivatives of \(\omega\). Cauchy--Schwarz gives bounds
\(\sqrt{M_4^*}/q_{\min}\) and \(2\sqrt{M_4^*}/q_{\min}\), respectively,
because \(\mathbb E(v^\top Z)^2=1\) and \(\mathbb E(|Z|^2-2)^2=4\).
Integrate along line segments in the parameter rectangle and average
clusters. The stated coefficient for \(\rho\) is conservative, since
only the second cluster mean changes. The fourth moment follows by
expanding \(|c+qZ|^4\) and using Gaussian second and fourth moments.
\(\square\)

Initial-mass variation is covered by \(\Delta(0)\): for aligned constant
masses, it equals \(|S_0-\widetilde S_0|\) when the initial angles are
identical.

\section{11. Proposition P9: probe rates require a gate-strip
modulus}\label{proposition-p9-probe-rates-require-a-gate-strip-modulus}

\subsection{11.1 Why state stability is not
enough}\label{why-state-stability-is-not-enough}

For a nonzero probe define

\[
 \Gamma_\nu(\xi,\eta)=
 \int\mathbf1_{\{|a_\theta^\top\xi|\le\eta|\xi|\}}\,d\nu,
 \qquad 0\le\eta\le1.
 \tag{11.1}
\]

For \(\eta>1\) set \(\Gamma_\nu(\xi,\eta)=S\). This measures
squared-radius mass near the probe's two gate boundaries, not Gaussian
training mass near a data hyperplane.

Small changes in angles make \(f(\xi)\) change continuously, but they
can change the indicator in the second term of (8.3). Therefore an
output or state enclosure alone does not automatically give a small
signed-rate error.

A concrete static example makes the obstruction explicit. Take
\(\xi=e_1\), \(\nu_n=M\delta_{(\pi/2-1/n,0)}\), and
\(\nu=M\delta_{(\pi/2,0)}\). Their outputs converge uniformly on the
unit circle. At the limit, the zero-gate convention gives
\(D_{1,\nu}(e_1)=0\). For \(P_1=\mathcal N(e_1,q^2I_2)\), however,

\[
 \lim_{n\to\infty}D_{1,\nu_n}(e_1)
 =-\frac M2\left(\frac{1+q^2}{2}-\frac{Mq}{\sqrt{2\pi}}\right),
\]

which is nonzero for sufficiently small \(M>0\). To calculate it, the
smooth first term of (8.3) vanishes in the limit, while the second tends
to \(M e_1 e_1^\top R_{1,\nu}(\pi/2)e_1\). Independence of \(X_1\) and
\(X_2\) gives the displayed bracket. This is a counterexample to generic
rate continuity at probe-boundary atoms, not a counterexample to
initialized Theorem A or a claim that its trajectory reaches this state.

\subsection{11.2 Explicit static rate perturbation
bound}\label{explicit-static-rate-perturbation-bound}

Use the quantities \(A_\theta,A_*,B_*,\eta_f\) from (10.1), and assume
the same data law and \(S,\widetilde S\le M\). Put

\[
 H_p=M_{2,p}(1+M),\qquad J_p=2W_p(1+M),
\]

\[
 \mathcal B_p=2H_pB_*
 +2M\bigl(M_{2,p}\eta_f+J_pA_\theta+H_pA_*\bigr)
 +H_p\Gamma_{\widetilde\nu}(\xi,A_\theta).
 \tag{11.2}
\]

Then

\[
 |V_{p,\nu}(\xi)-V_{p,\widetilde\nu}(\xi)|
 \le |\xi|\mathcal B_p,
 \tag{11.3}
\]

\[
 \boxed{
 |D_{p,\nu}(\xi)-D_{p,\widetilde\nu}(\xi)|
 \le |\xi|^2\left[
 MH_p\eta_f+\frac{1+M}{2}\mathcal B_p\right].}
 \tag{11.4}
\]

\textbf{Proof.} Use (8.3). Its unit-mass integrand has norm at most
\(2H_p|\xi|\), which bounds the mass-difference contribution by
\(2H_pB_*|\xi|\). At matched labels, the smooth difference of the first
term is at most

\[
 |\xi|\left(\|R_p-\widetilde R_p\|+2H_pA_\theta\right).
\]

For the second term, first hold its probe indicator fixed. Since
\(\|bb^\top-\widetilde b\widetilde b^\top\|\le2|b-\widetilde b|\), the
smooth difference is at most

\[
 |\xi|\left(\|R_p-\widetilde R_p\|+2H_pA_\psi\right).
\]

If the two probe indicators differ, then
\(|a_{\widetilde\theta}^\top\xi|\le A_\theta|\xi|\). The remaining
integrand is at most \(H_p|\xi|\) there. Integrate it with the reference
mass to get the last term in (11.2). Finally use the cluster version of
(10.5) and \(\widetilde S\le M\) to prove (11.3).

For (11.4), add and subtract the product
\(\widetilde e(\xi)^\top V_{p,\nu}(\xi)\) in (8.4). Use
\(|e-\widetilde e|\le\eta_f|\xi|\), \(|V_{p,\nu}|\le2MH_p|\xi|\), and
\(|\widetilde e|\le(1+M)|\xi|\). \(\square\)

For different Gaussian laws, add
\((1+M)\|\omega_p-\widetilde\omega_p\|_1\) to the residual-matrix bound
inside (11.2), and use constants that bound both laws. Thus parameter
openness of strict rate signs follows when the reference strip modulus
is controlled; it is not obtained from (10.6) alone.

\subsection{11.3 Sector variation}\label{sector-variation}

For a fixed state with mass \(S\le M\) and unit probes
\(\xi,\widetilde\xi\), let \(\varepsilon=|\xi-\widetilde\xi|\). The same
splitting gives

\[
 |V_{p,\nu}(\xi)-V_{p,\nu}(\widetilde\xi)|
 \le 2MH_p\varepsilon+H_p\Gamma_\nu(\widetilde\xi,\varepsilon).
 \tag{11.5}
\]

Here a changed indicator implies
\(|a_\theta^\top\widetilde\xi|\le\varepsilon\). Since \(e_\nu\) has
Lipschitz constant at most \(1+M\),

\[
 |D_{p,\nu}(\xi)-D_{p,\nu}(\widetilde\xi)|
 \le2MH_p(1+M)\varepsilon
 +\frac12(1+M)H_p\Gamma_\nu(\widetilde\xi,\varepsilon).
 \tag{11.6}
\]

Uniform strip bounds therefore give uniform sector continuity. An error
curve is easier to control: for unit probes,
\(|E_\xi-E_{\widetilde\xi}|\le(1+S)^2|\xi-\widetilde\xi|\) directly,
without a gate-strip term.

\subsection{11.4 Initialization and a short-time strip
bound}\label{initialization-and-a-short-time-strip-bound}

At aligned isotropic initialization,

\[
 \Gamma_{\nu_0}(\xi,\eta)
 =\frac{2S_0}{\pi}\arcsin\eta,\qquad0\le\eta\le1,
 \tag{11.7}
\]

uniformly over nonzero probes. This is because the strip consists of two
arcs of total angular length \(4\arcsin\eta\).

A short-time extension can also be proved without assuming it. Fix
\(T>0\), use

\[
 M=S_0+M_2T/2,\quad H=M_2(1+M),\quad J=2W(1+M),
 \quad C_{\rm ang}=2(H+J).
\]

The angular vector field, with the already-determined measure path held
fixed, is \(C^1\) in \((\theta,\psi)\). Its Jacobian has induced
\(\ell^1\) norm at most \(C_{\rm ang}\): the derivatives in \(\theta\)
have magnitude at most \(H+J\) in each equation, and derivatives in
\(\psi\) at most \(H\) in each equation. For the initial-label
derivative \(Z=\partial_\alpha(\theta,\psi)\),

\[
 |Z(t)|_1\le2e^{C_{\rm ang}t},\qquad
 |\partial_\alpha\theta_t-1|\le e^{C_{\rm ang}t}-1.
 \tag{11.8}
\]

Indeed \(\dot Z=Dv_t Z\), \(Z(0)=(1,1)\); integration of the first-row
bound gives the second inequality. Thus on

\[
 0\le t\le T_{\rm short}:=
 \min\left\{T,\frac{\log(3/2)}{C_{\rm ang}}\right\},
\]

one has \(\partial_\alpha\theta_t\ge1/2\). The map
\(\alpha\mapsto\theta_t(\alpha)\) is then a degree-one circle
diffeomorphism. Its input-angle mass density is at most
\(S_0e^{2C_0T}/\pi\), giving

\[
 \Gamma_{\nu_t}(\xi,\eta)
 \le\min\left\{S(t),\frac{4S_0e^{2C_0T}}{\pi}\arcsin\eta\right\}
 \quad(0\le t\le T_{\rm short}).
 \tag{11.9}
\]

More generally, an independently established input-angle mass-density
bound \(B_T\) gives \(\Gamma\le4B_T\arcsin\eta\). The joint angular flow
being invertible does not by itself supply such a bound for its
projection at late times. Proving a useful late strip bound or enclosing
gate-boundary contributions directly remains part of the
reachability/certification work.

\section{12. Scope of the original P1--P9
layer}\label{scope-of-the-original-p1p9-layer}

Sections 12--14 describe the original foundations layer. Sections 15--24
below extend that layer; Section 24 gives the updated proof status.

The derivations provide a locked model, exact finite-to-continuum
normalization, aligned initial values, the full angular
transport--reaction equations, exact Gaussian integral formulas, global
well-posedness in the specified characteristic class, loss dissipation,
finite-horizon mass and speed bounds, exact cluster rates and learning
observables, labeled-state/data stability, and a gate-aware
rate-stability estimate. They also provide an explicit early-time strip
bound.

They do not establish a canonical learned-state tube, a positive-time
strong-progress lower bound, substantial weak learning, the persistence
of \(D_1<0<D_2\) into a learned regime, dominance crossover, a useful
late gate-strip bound, or a quantitative reversal amplitude. They do not
establish Theorem B, a normal-form timing theorem, or a
finite-sample/finite-width high-probability lifting theorem.

In particular, \textbf{a population pilot that integrates these
equations is a witness search, not a proof}. A validated pilot would
additionally need rigorous angular quadrature error, time-integration
defect, continuation control, and probe-rate/gate enclosures. The crude
global Gronwall constants here are not claimed to deliver such a
certificate at the canonical parameters.

\section{13. Immediate next proof
interface}\label{immediate-next-proof-interface}

The next computation should use exactly (4.4), with (6.8)--(6.9) or
(6.11), and evaluate (8.4), (9.1), and (9.2). The first audit should
check mass normalization, the \(S_0/4\) initial map, the factor \(1/2\)
in each cluster contribution, and agreement of the matrix-kernel and
direct-velocity formulas.

A subsequent reachability proof must output a state enclosure and bounds
on the terms in (11.4), not just a visually close output trajectory.
Only after that enclosure is available should one certify
common-interval opposite signs, early/late dominance, and the
weak-learning floor.

The proof order is therefore:

\[
 \begin{gathered}
 \text{locked setup and exact equations}\\
 \Downarrow\\
 \text{well-posed population flow and gate-aware stability}\\
 \Downarrow\\
 \text{population witness and initialized reachability}\\
 \Downarrow\\
 \text{Theorem A}.
 \end{gathered}
\]

The foundations are not a substitute for the central dynamical argument.
Their role is to make every quantity in that argument exact and to
expose, rather than hide, the remaining proof obligations.

\section{14. Non-certifying algebra
check}\label{non-certifying-algebra-check}

A separate floating-point check compared the identities at a synthetic
seven-neuron balanced state under positive Gaussian noise. It used
one-dimensional angular quadrature split at the training gate
boundaries, and compared the direct Cartesian gradient, angular
equations, closed Gaussian \(B_p\), the \(B_p,C_p\) closure, both
cluster-velocity formulas, the mass identity, and centered finite
differences. The direct algebraic comparisons agreed within
approximately \(3\cdot10^{-15}\); the finite-difference comparisons
within approximately \(10^{-9}\) at the tested step sizes.

This check is a transcription/normalization diagnostic only. It is not
interval arithmetic, does not establish uniform bounds, does not replace
the proofs, and did not integrate an initialized population trajectory
or test the crossover conjecture.

\section{15. Supplement to P6: exact gradient attribution and the
correct
metric}\label{supplement-to-p6-exact-gradient-attribution-and-the-correct-metric}

The results in Sections 3--11 are retained. The additions below concern
exact identities, controls, and conditional certificate interfaces; they
do not assert that the canonical initialized trajectory has reached the
competition regime. In particular, the global constants of P8 and the
short-time strip estimate of P9 are not proposed as a quantitative
continuation certificate over the canonical reversal horizon.

Let \(Z=(U,W)\) denote the labeled Cartesian state, and use the
fixed-reference mean-field inner product

\[
 \langle (v,z),(\widetilde v,\widetilde z)\rangle_{\mathcal H}
 =\int(v^\top\widetilde v+z^\top\widetilde z)\,d\lambda.
 \tag{15.1}
\]

At states without label mass on the probe boundary, the first variations
of \(E_\xi\) are represented by

\[
 (\nabla_{\mathcal H}E_\xi)_U
 =-\xi(W^\top e(\xi))\mathbf1_{\{U^\top\xi>0\}},
 \qquad
 (\nabla_{\mathcal H}E_\xi)_W=-e(\xi)(U^\top\xi)_+.
 \tag{15.2}
\]

The single-cluster loss \(\mathcal L_p=\frac12\mathbb E_{P_p}|e(X)|^2\)
has gradient

\[
 (\nabla_{\mathcal H}\mathcal L_p)_U=-R_p^\top W,
 \qquad
 (\nabla_{\mathcal H}\mathcal L_p)_W=-R_pU.
 \tag{15.3}
\]

Consequently,

\[
 \boxed{D_p=-\frac12\langle\nabla_{\mathcal H}E_\xi,
                         \nabla_{\mathcal H}\mathcal L_p\rangle_{\mathcal H}.}
 \tag{15.4}
\]

\textbf{Proof.} Differentiate \(f_Z(\xi)=\int W(U^\top\xi)_+d\lambda\)
against bounded variations. The zero-boundary-mass assumption and
dominated convergence justify the first variation, giving (15.2).
Applying the same calculation inside the Gaussian training expectation
gives (15.3). Substitution into (15.1) gives
\(\langle\nabla E,\nabla\mathcal L_p\rangle=e(\xi)^\top V_p(\xi)\), with
\(V_p\) exactly as in (8.3). Equation (8.4) then yields (15.4).
Equivalently, \(\dot Z=-\frac12(\nabla\mathcal L_1+\nabla\mathcal L_2)\)
and the chain rule gives the same decomposition. \(\square\)

Thus \(D_2>0\) means that the weak-cluster descent direction increases
probe error instantaneously at the current full state. It is a negative
inner product of gradients, not a counterfactual statement about
retraining on cluster 2 alone. At a probe-boundary atom, use a specified
compatible selected first variation and the almost-everywhere pathwise
convention of P6 rather than assert an ordinary differentiable gradient
at that state.

On the balanced parameter manifold, the Cartesian metric pulls back to

\[
 \|\delta Z\|_{\mathcal H}^2
 =\int\left[\frac{(\delta m)^2}{2m}
       +m(\delta\theta)^2+m(\delta\psi)^2\right]d\lambda.
 \tag{15.5}
\]

Indeed, differentiating \(U=\sqrt m\,a_\theta\) and
\(W=\sqrt m\,b_\psi\) gives two radial squared contributions
\((\delta m)^2/(4m)\) and the two tangential contributions in (15.5). A
Euclidean logarithmic norm of the \((\theta,\psi,m)\) ODE is therefore
not automatically the relevant gradient-flow one-sided bound. A
time-dependent polar metric also contributes its own metric derivative
to a stability calculation.

\section{16. Proposition P10: exact matched linear
control}\label{proposition-p10-exact-matched-linear-control}

Replace ReLU by identity activation, keeping the untied architecture,
the population data law, and the aligned isotropic initialization. In
Cartesian mean-field coordinates, let

\[
 M=\int WU^\top d\lambda,\qquad f(x)=Mx,\qquad
 A=\operatorname{diag}(a_1,a_2),
\]

\[
 a_1=\tfrac12+q^2,\qquad a_2=\tfrac{\rho^2}{2}+q^2.
\]

For any \(S_0>0\) the initialized solution remains aligned, and

\[
 M(t)=\operatorname{diag}(\ell_1(t),\ell_2(t)),\qquad
 \ell_k(0)=S_0/2,
\]

\[
 \dot\ell_k=2a_k\ell_k(1-\ell_k),\qquad
 \ell_k(t)=\frac{1}{1+(2/S_0-1)e^{-2a_kt}}.
 \tag{16.1}
\]

Every fixed-probe error is non-increasing:

\[
 E_\xi(t)=\frac12\sum_{k=1}^2(1-\ell_k(t))^2\xi_k^2,
\]

\[
 \boxed{\dot E_\xi(t)=-2\sum_{k=1}^2a_k\ell_k(t)(1-\ell_k(t))^2\xi_k^2\le0.}
 \tag{16.2}
\]

There is a stronger clusterwise conclusion. Put

\[
 A_1=\operatorname{diag}(1+q^2,q^2),\qquad
 A_2=\operatorname{diag}(q^2,\rho^2+q^2).
\]

Then, with the same cluster weight as in (15.4),

\[
 \boxed{D_p^{\rm lin}(t,\xi)
 =-\sum_{k=1}^2(A_p)_{kk}\ell_k(t)(1-\ell_k(t))^2\xi_k^2\le0.}
 \tag{16.3}
\]

\textbf{Proof.} The exact Cartesian field is \(\dot W=RU\),
\(\dot U=R^\top W\), with \(R=(I-M)A\). Construct an aligned solution
\(U(t,\alpha)=W(t,\alpha)=T(t)U(0,\alpha)\), where \(T\) is diagonal and
\(T_{kk}=\sqrt{2\ell_k/S_0}\). Isotropic initialization gives
\(\int U(0)U(0)^\top d\lambda=S_0I_2/2\), hence its moment is
\(M=\operatorname{diag}(\ell_k)\). Equation (16.1) implies
\(\dot T_{kk}=a_k(1-\ell_k)T_{kk}\), exactly the Cartesian field because
\(R\) is diagonal. The explicit denominator in (16.1) is positive for
every \(t\ge0\) and \(S_0>0\). Local uniqueness of the Cartesian field
identifies this constructed solution with the initialized flow.
Differentiating \(E_\xi\) proves (16.2).

For a raw single-cluster field, \(R_p=(I-M)A_p\) is diagonal. Its
contribution to \(\dot M\), before applying the mixture weight, is
\(R_pM+MR_p=2\operatorname{diag}((A_p)_{kk}\ell_k(1-\ell_k))\).
Multiplying the resulting probe velocity by the cluster weight \(1/2\)
and contracting with \(-e(\xi)\) yields (16.3). \(\square\)

\textbf{Normalization.} With the same parameter initialization mass, the
linear initial map is \(S_0x/2\), whereas the ReLU initial map is
\(S_0x/4\). An initial-output-matched linear control uses linear mass
\(S_0/2\). Both matching conventions yield the monotonicity conclusion.
For \(S_0>2\), linear coordinates approach 1 from above rather than
below; (16.2) still holds. At \(S_0=2\) the linear model is initially
the identity.

\subsection{16.1 Do not extend the diagonal closure off the commuting
manifold}\label{do-not-extend-the-diagonal-closure-off-the-commuting-manifold}

For a general untied linear state set

\[
 P_U=\int UU^\top d\lambda,\qquad P_W=\int WW^\top d\lambda.
\]

The actual equation is

\[
 \boxed{\dot M=(I-M)AP_U+P_W(I-M)A.}
 \tag{16.4}
\]

This follows by differentiating \(WU^\top\): \(\dot M=RP_U+P_WR\). It is
not a closed equation in \(M\) alone. Even at an instant when \(U=W\)
and \(M\) is symmetric, the resulting derivative is

\[
 (I-M)AM+M(I-M)A,
\]

and alignment generally fails immediately if the covariance is full rank
and \([A,M]\ne0\). Indeed \(\dot W-\dot U=(R-R^\top)U\), so full-rank
aligned covariance and tangency require \(R=R^\top\), equivalently
\([A,M]=0\).

Along the commuting aligned trajectory, (16.4) reduces to \(AM+MA-2MAM\)
and to (16.1). The commuting class includes any diagonal \(M\) in the
eigenbasis of \(A\), not only scalar-isotropic \(M\).

For comparison only, the polynomial
\(F_{\rm SIM}(M)=AM+MA-\frac12(AM^2+M^2A+2MAM)\) printed in {[}S1, Eq.
(D.2){]} obeys

\[
 (AM+MA-2MAM)-F_{\rm SIM}(M)=\tfrac12[[A,M],M].
 \tag{16.5}
\]

That is an algebraic identity between polynomials, not a derivation of
an off-commuting untied flow. Direct differentiation of
\(\frac12\|(UU^\top-I)A^{1/2}\|_F^2\) in tied Euclidean parameters gives
twice \(F_{\rm SIM}\) as its continuous-time effective-matrix field.
None of P10 depends on importing that normalization.

\section{17. Proposition P10b: ReLU creates a probe-active minor entry
immediately}\label{proposition-p10b-relu-creates-a-probe-active-minor-entry-immediately}

For a nonzero probe define the exact active matrix

\[
 M_\xi(t)=\int\mathbf1_{\{a_\theta^\top\xi>0\}}b_\psi a_\theta^\top d\nu_t,
 \qquad f_t(\xi)=M_\xi(t)\xi.
 \tag{17.1}
\]

At aligned isotropic initialization,

\[
 M_\xi(0)=\frac{S_0}{4}I_2\quad\text{for every nonzero }\xi.
 \tag{17.2}
\]

\textbf{Proof.} The map \(a\mapsto aa^\top\) is unchanged by
\(a\mapsto-a\), and exactly one point of each antipodal pair lies in the
active half-circle, except on a zero-measure boundary. The integral on
any active half-circle is therefore half the full second moment
\(S_0I_2/2\). \(\square\)

In the normalized two-concept Gaussian model, for \(q>0\), \(\rho>0\),
and \(0<S_0<4\),

\[
 \boxed{\left.\frac d{dt}[M_{e_1}(t)]_{21}\right|_{t=0}
       =\left.\frac d{dt}f_{t,2}(e_1)\right|_{t=0}>0.}
 \tag{17.3}
\]

Thus a nonzero probe-active off-diagonal response is generated in the
exact continuum flow without initial random off-diagonal fluctuations.
This is an initial generation result, not a reversal theorem.

\textbf{Proof.} Write \(c_0=1-S_0/4>0\). Initially
\(R(\theta)=c_0\mathbb E[XX^\top\mathbf1_{a_\theta^\top X>0}]\) is
symmetric, so the initial velocities of \(U\) and \(W\) agree. Define

\[
 H(x)=\frac1{2\pi}\int_{-\pi/2}^{\pi/2}
 (\cos\alpha\,x_2+\sin\alpha\,x_1)
 (a_\alpha^\top x)_+\,d\alpha.
 \tag{17.4}
\]

Differentiating \(f_2(e_1)\) in the two parameter layers gives

\[
 \dot f_{0,2}(e_1)=c_0S_0\mathbb E_P H(X).
 \tag{17.5}
\]

There is no probe-boundary atom at initialization. The preceding
derivative is therefore justified by the first-variation calculation of
Section 15.

For \(x=r(\cos\beta,\sin\beta)\) with \(0\le\beta\le\pi\), the active
integration interval is \((\beta-\pi/2,\pi/2)\). The integrand becomes
\(r^2\sin(\alpha+\beta)\cos(\alpha-\beta)\). Using
\(\sin A\cos B=[\sin(A+B)+\sin(A-B)]/2\) and integrating yields

\[
 H(x)=\frac{r^2}{4\pi}
       [\sin^2\beta+(\pi-\beta)\sin(2\beta)].
 \tag{17.6}
\]

Reflection of \(\alpha\) in (17.4) shows \(H(x_1,-x_2)=-H(x_1,x_2)\). If
\(x_1\ge0\) and \(x_2>0\), write \(0<\beta\le\pi/2\). Then

\[
 H(x_1,x_2)+H(-x_1,x_2)
 =\frac{r^2}{2\pi}
 [\sin^2\beta+(\pi-2\beta)\sin\beta\cos\beta]>0.
 \tag{17.7}
\]

Cluster 1 is symmetric in \(X_2\), so \(\mathbb E_{P_1}H=0\). Cluster 2
is symmetric in \(X_1\) and, for \(x_2>0\), its density satisfies
\(p_2(x_1,x_2)>p_2(x_1,-x_2)\) because \(\rho>0\). Integrating over
\(x_1\ge0,x_2>0\), using (17.7) and oddness in \(x_2\), gives
\(\mathbb E_{P_2}H>0\). Equation (17.5) now proves (17.3). Finally
\([M_{e_1}]_{21}=f_2(e_1)\) identically. \(\square\)

As a normalization check, in the noiseless two-point limit the same
calculation gives \(\dot f_{0,2}(e_1)=c_0S_0\rho^2/(8\pi)\). The
positive-noise claim above does not use a small-\(q\) perturbation.

\textbf{Vocabulary and flux.} It is consistent to reserve ``major'' for
diagonal entries and ``minor'' for off-diagonal entries of \(M_\xi\). A
transverse diagonal entry of a strong cohort block is not, for that
reason alone, a minor entry of \(M_\xi\). Nor does (17.2) make all
training-and-probe doubly gated moments diagonal.

Although \(f_t(\xi)=M_\xi(t)\xi\) is exact, differentiating \(M_\xi\)
itself can create a boundary-flux matrix. Formally it contains
\(\int WU^\top\delta(U^\top\xi)(\dot U^\top\xi)d\lambda\). Its
contraction with \(\xi\) vanishes because \(U^\top\xi=0\) on the
boundary. Thus there is no additional output-rate flux term in P6, but
matrix-entry equations still require the boundary contribution or a
proved no-flux condition.

\section{18. Proposition P11: explicit Gaussian angular tails and
interior gate
purity}\label{proposition-p11-explicit-gaussian-angular-tails-and-interior-gate-purity}

\subsection{18.1 Angular tails, including the backward
half-plane}\label{angular-tails-including-the-backward-half-plane}

Let \(c\ne0\), \(c_*=|c|\), and \(\gamma\) be the angular distance from
\(a_\beta\) to \(c/c_*\). Put \(\ell=c_*\cos\gamma\). For \(\ell\ge0\),
(6.3) gives

\[
 \omega_c(\beta)\le W_c
 \exp\left[-\frac{c_*^2\sin^2\gamma}{2q^2}\right].
 \tag{18.1}
\]

For \(\ell\le0\) there is a useful stronger backward bound:

\[
 \omega_c(\beta)\le\frac{q^2}{\pi}
                 e^{-c_*^2/(2q^2)}.
 \tag{18.2}
\]

\textbf{Proof.} Integration by parts gives
\(\partial_\ell J_3=3J_2\ge0\), so \(J_3(\ell,q)\le J_3(c_*,q)\) when
\(0\le\ell\le c_*\). This proves (18.1). For \(\ell\le0\),

\[
 (r-\ell)^2=r^2+\ell^2-2r\ell\ge r^2+\ell^2,
\]

so \(J_3(\ell,q)\le e^{-\ell^2/(2q^2)}J_3(0,q)=2q^4e^{-\ell^2/(2q^2)}\).
Substitute into (6.3) to prove (18.2). \(\square\)

For an angular set \(B\) outside the cone \(\gamma<\delta\), with
\(0<\delta\le\pi/2\), it follows that

\[
 \int_B\omega_c(\beta)d\beta
 \le |B|\max\left\{W_c e^{-c_*^2\sin^2\delta/(2q^2)},
                    (q^2/\pi)e^{-c_*^2/(2q^2)}\right\}.
 \tag{18.3}
\]

Here \(|B|\) is its Lebesgue angular length. The separate backward
estimate avoids mistaking the antipodal direction for a second
concentration peak.

\subsection{18.2 Exact wrong-gate second
moments}\label{exact-wrong-gate-second-moments}

For \(X\sim\mathcal N(c,q^2I_2)\), a unit gate vector \(a\), and
\(\ell=a^\top c\), set \(\kappa=\ell/q\). Then

\[
 \mathbb E[|X|^2\mathbf1_{\{a^\top X>0\}}]
 =(|c|^2+2q^2)\Phi(\kappa)+\ell q\phi(\kappa),
 \tag{18.4}
\]

\[
 \mathbb E[|X|^2\mathbf1_{\{a^\top X\le0\}}]
 =(|c|^2+2q^2)\Phi(-\kappa)-\ell q\phi(\kappa).
 \tag{18.5}
\]

\textbf{Proof.} Decompose \(X=Ya+Za^\perp\), where
\(Y\sim\mathcal N(\ell,q^2)\) and
\(Z\sim\mathcal N((a^\perp)^\top c,q^2)\) are independent. Direct
truncated-normal integration gives
\(\mathbb E[Y^2\mathbf1_{Y>0}]=(\ell^2+q^2)\Phi(\kappa)+\ell q\phi(\kappa)\).
Add \(\mathbb E[Z^2]\mathbb P(Y>0)\) to obtain (18.4); subtract from
\(\mathbb E|X|^2\) to obtain (18.5). \(\square\)

In particular, if \(\ell\le-b<0\), the wrong active-gate moment is at
most \((|c|^2+2q^2)\Phi(-b/q)\). If \(\ell\ge b>0\), the missing-gate
moment has the same bound.

Let \(0<\delta<\pi/4\) and consider the strict quadrant-II interior

\[
 J_\delta=[\pi/2+\delta,\pi-\delta].
\]

For every state of mass \(S\) and every \(\theta\in J_\delta\), define
the ungated weak residual moment
\(\overline R_2=\mathbb E_{P_2}[e(X)X^\top]\). Then

\[
 \boxed{\|R_1(\theta)\|\le
 (1+S)M_{2,1}\Phi(-\sin\delta/q),}
\]

\[
 \boxed{\|R_2(\theta)-\overline R_2\|\le
 (1+S)M_{2,2}\Phi(-\rho\sin\delta/q).}
 \tag{18.6}
\]

\textbf{Proof.} In \(J_\delta\),
\(a_\theta^\top c_1=\cos\theta\le-\sin\delta\) and
\(a_\theta^\top c_2=\rho\sin\theta\ge\rho\sin\delta\). Apply
(18.4)--(18.5) and \(|e(X)|\le(1+S)|X|\) to the relevant matrix
integrals. \(\square\)

This is a static gate-purity estimate for a strictly interior arc. It
does not claim purity uniformly up to \(\theta=\pi/2\), where the
strong-cluster gate probability is \(1/2\). It also does not prove that
sufficient initialized mass remains in this arc, acquires useful output
orientation, or yields an \(O(1)\) weak response. Those are the
dynamical obligations needed for weak learning.

\section{19. Exact shared-gate forcing and a signed rotation
interface}\label{exact-shared-gate-forcing-and-a-signed-rotation-interface}

Let \(\mathcal K(\kappa)=\phi(\kappa)+\kappa\Phi(\kappa)\), and set

\[
 \kappa=\frac{\rho\sin\theta}{q}.
\]

Equation (6.7) specializes to

\[
 \boxed{B_2(\theta)=\rho q\mathcal K(\kappa)e_2
                        +q^2\Phi(\kappa)a_\theta,}
\]

\[
 G_{2,\nu}(\theta)=B_2(\theta)
             -\int b_\varphi C_2(\theta,\vartheta)d\nu(\vartheta,\varphi).
 \tag{19.1}
\]

This keeps the gate-overlap target forcing and the fitted-output
subtraction distinct. At \(\kappa=0.234\),
\(\mathcal K(\kappa)\simeq0.527\); this numerical substitution is not a
dynamically proved tilt.

The cluster-2 contribution to output-angle velocity is exactly

\[
 (\dot\psi)_2=\tfrac12(b_\psi^\perp)^\top G_{2,\nu}(\theta).
 \tag{19.2}
\]

A positive \(e_2\) component of \(B_2\) alone does not determine the
sign of (19.2), or of \(D_2\). For a scalar \(r\ge0\) define

\[
 \varepsilon_2(\theta;r)=
 \mathbb E_{P_2}[(e(X)-rX_2e_2)(a_\theta^\top X)_+],
\]

\[
 T_2(\theta)=\mathbb E_{P_2}[X_2(a_\theta^\top X)_+]
 =\rho q\mathcal K(\kappa)+q^2\sin\theta\Phi(\kappa).
\]

Then the exact identity and a sufficient positive-rotation condition are

\[
 G_2(\theta)=rT_2(\theta)e_2+\varepsilon_2(\theta;r),
\]

\[
 (\dot\psi)_2\ge\tfrac12
 [rT_2(\theta)(b_\psi^\perp)^\top e_2-|\varepsilon_2(\theta;r)|],
 \tag{19.3}
\]

with the explicit bound

\[
 |\varepsilon_2(\theta;r)|
 \le \sqrt{C_2(\theta,\theta)}
      \|e-rX_2e_2\|_{L^2(P_2)}.
 \tag{19.4}
\]

The proof is substitution and Cauchy--Schwarz. A bound on the scalar
correlation statistic \(\mathsf m_2\) alone does not imply the residual
approximation in (19.4). A signed mechanism proof must control that
approximation or retain the exact selected residual integral.

\section{20. Proposition P12: signed residual-curvature identity and
one-sided
continuation}\label{proposition-p12-signed-residual-curvature-identity-and-one-sided-continuation}

\subsection{20.1 Exact second variation in Cartesian
coordinates}\label{exact-second-variation-in-cartesian-coordinates}

Let \(Z=(U,W)\) be a continuous labeled Cartesian state, not necessarily
balanced, with \(\inf_\alpha|U(\alpha)|>0\) and bounded \(U,W\). Let
\(h=(v,z)\) be a bounded continuous perturbation. All inner products
below use (15.1); the data law is fixed. Put
\(F(Z)=-\nabla_{\mathcal H}\mathcal L(Z)\) and define its output
differential

\[
 J_Zh(x)=\int\left[z(U^\top x)_+
             +W\mathbf1_{\{U^\top x>0\}}(v^\top x)\right]d\lambda.
 \tag{20.1}
\]

Write \(r=|U|\), \(a=U/r=a_\theta\), \(n=a_\theta^\perp\). Define

\[
 h_Z^{\rm gate}(\alpha)=
 \sum_{\epsilon\in\{-1,1\}}
 \omega(\theta+\epsilon\pi/2)
 \frac{W(\alpha)^\top e_Z(a_{\theta+\epsilon\pi/2})}{r(\alpha)}.
 \tag{20.2}
\]

Then

\[
 \boxed{\begin{aligned}
 \langle h,DF(Z)h\rangle_{\mathcal H}
 ={}&-\|J_Zh\|_{L^2(P)}^2
       +2\int z^\top R_Z(\theta)v\,d\lambda\\
    &+\int h_Z^{\rm gate}(\alpha)(v^\top n)^2d\lambda.
 \end{aligned}}
 \tag{20.3}
\]

The signs in (20.3) use \(e=x-f\), as in the rest of this document. Thus
the Gauss--Newton term is nonpositive, while the residual terms need not
have either sign.

\textbf{Proof.} For \(Z_s=Z+sh\), the first output variation is (20.1).
The second output variation, interpreted inside the Gaussian
expectation, is

\[
 \delta^2f(x)=\int\left[
 2z\mathbf1_{\{U^\top x>0\}}(v^\top x)
 +W\delta(U^\top x)(v^\top x)^2\right]d\lambda.
\]

This distributional identity may be justified without assigning any
finite mass to an individual training hyperplane: use polar coordinates
and differentiate the moving half-circle endpoints, exactly as in
(6.12). Equivalently, the degree-two feature kernels of P4 are \(C^2\)
for nonzero input vectors. Their derivatives justify differentiation of
the loss twice on a uniformly nonzero, bounded Cartesian tube. Dominated
integration over labels is valid there.

Since \(e=x-f\), the negative Hessian quadratic form is

\[
 -\delta^2\mathcal L=-\mathbb E|J_Zh|^2+
                          \mathbb E[e^\top\delta^2f].
\]

The smooth bilinear part is \(2\int z^\top R_Zv\,d\lambda\). For the
boundary part write \(X=s a_\beta\), with \(s>0\). The delta factor is
\(\delta(U^\top X)=(rs)^{-1}\delta(a_\theta^\top a_\beta)\). Homogeneity
of \(e\), the two perturbation factors, and the polar Jacobian leave the
radial integral \(\omega(\beta)\). The two roots are
\(\beta=\theta\pm\pi/2\), each with unit angular derivative magnitude,
and \((v^\top a_\beta)^2=(v^\top n)^2\). The remaining coefficient is
precisely (20.2). This proves (20.3). \(\square\)

The quadratic form is bounded and extends to the corresponding
\(L^2(\lambda)\) perturbations on such a tube. Differentiability is
asserted in the continuous bounded-state setting used to derive the
formula, not on an unrestricted \(L^2\)-open set that would allow
vanishing input weights on small label subsets.

For a balanced state, \(W/r=b_\psi\), so

\[
 h_Z^{\rm gate}=
 \sum_{\epsilon=\pm1}\omega(\theta+\epsilon\pi/2)
                    b_\psi^\top e_Z(a_{\theta+\epsilon\pi/2}).
 \tag{20.4}
\]

This is a \textbf{signed projection} of the boundary residual. Large
\(\omega\) alone is not a positive one-sided expansion rate. In
particular, replacing (20.4) by a norm of the full boundary residual can
discard both a favorable sign and near-orthogonality to the output
direction.

\textbf{Exact initial cancellation.} At the aligned isotropic
initialization, \(W/r=a_\theta\) and \(e_0(x)=c_0x\), with
\(c_0=1-S_0/4\). Every boundary direction is perpendicular to
\(a_\theta\), so

\[
 h_{Z_0}^{\rm gate}(\alpha)=0\quad\text{for every label}.
 \tag{20.4a}
\]

For \(0<S_0<4\), the gated initial moment is
\(R_0(\theta)=c_0\mathbb E[XX^\top\mathbf1_{a_\theta^\top X>0}]\), which
is positive semidefinite and bounded above by \(c_0A\). Using
\(2z^\top R_0v\le\|R_0\|(|z|^2+|v|^2)\) and retaining or discarding the
nonpositive Gauss--Newton term gives the initial bound

\[
 \langle h,DF(Z_0)h\rangle_{\mathcal H}
 \le c_0\|A\|\|h\|_{\mathcal H}^2
 =c_0(1/2+q^2)\|h\|_{\mathcal H}^2.
 \tag{20.4b}
\]

Thus a large \(O(1/q)\) norm of \(\partial_\theta R\) does not imply an
\(O(1/q)\) expanding one-sided rate at initialization. This cancellation
is exact at the specified initial state; retaining it approximately on a
nonzero tube requires further bounds.

\subsection{20.2 A local operator suitable for
certification}\label{a-local-operator-suitable-for-certification}

The residual part of (20.3) is the multiplication quadratic form
associated with

\[
 \mathcal A_Z(\alpha)=
 \begin{pmatrix}
 h_Z^{\rm gate}nn^\top&R_Z^\top\\
 R_Z&0
 \end{pmatrix}
 \quad\text{on }(v,z).
 \tag{20.5}
\]

A valid upper one-sided rate is
\(\operatorname*{ess\,sup}_\alpha\lambda_{\max}(\mathcal A_Z(\alpha))\).
A potentially sharper rate retains the nonlocal negative term:

\[
 \ell_Z=\sup_{\|h\|_{\mathcal H}=1}
 \left\{-\|J_Zh\|_{L^2(P)}^2
      +\int h(\alpha)^\top\mathcal A_Z(\alpha)h(\alpha)d\lambda\right\}.
 \tag{20.6}
\]

These are exact/sufficient interfaces. No claim is made that their upper
bounds are small or negative along the canonical trajectory. Any
numerical upper bound must enclose the continuum integrals, not only
eigenvalues of a point-sampled finite matrix.

\subsection{20.3 Conditional one-sided tube
lemma}\label{conditional-one-sided-tube-lemma}

Let \(\widehat Z(t)\) be a differentiable reference and
\(d(t)=\dot{\widehat Z}(t)-F(\widehat Z(t))\) its defect. Suppose a
prescribed labelwise tube contains every segment joining
\(\widehat Z(t)\) to the candidate true state, keeps all input radii
bounded away from zero, and satisfies

\[
 \langle h,DF(Y)h\rangle_{\mathcal H}
 \le\ell(t,R)\|h\|_{\mathcal H}^2
 \tag{20.7}
\]

for all states \(Y\) on those segments when the error norm is at most
\(R\). If \(r(t)=\|Z(t)-\widehat Z(t)\|_{\mathcal H}\), then

\[
 D^+r(t)\le\ell(t,r(t))r(t)+\|d(t)\|_{\mathcal H}.
 \tag{20.8}
\]

\textbf{Proof.} Put \(h=Z-\widehat Z\) and integrate \(DF\) along the
Cartesian parameter segment. Then

\[
 \tfrac12\frac d{dt}\|h\|_{\mathcal H}^2
 =\langle h,F(Z)-F(\widehat Z)\rangle-\langle h,d\rangle
 \le\ell(t,r)r^2+r\|d\|.
\]

Divide by \(r\) when \(r>0\); at \(r=0\) the same inequality follows for
the upper right derivative of the norm. \(\square\)

A strict supersolution of (20.8), together with proved labelwise side
guards, closes the usual first-exit comparison. An \(L^2\) radius alone
does not guarantee a pointwise positive input radius, and does not
provide P9's uniform angular error budget. Those labelwise guards and a
compatible probe-rate error bound are additional parts of a certificate.
Cartesian line segments between balanced states generally are not
balanced, which is why (20.2), not only (20.4), is required in a tube
proof.

\section{21. Proposition P13: projected density, covariance, and probe
strips}\label{proposition-p13-projected-density-covariance-and-probe-strips}

Let \(\sigma_t\) denote the input-angle marginal of \(\nu_t\). For a
regular angle \(\vartheta\), if the preimages of \(\vartheta\) under
\(\alpha\mapsto\theta_t(\alpha)\) are noncritical, the local density is

\[
 n_t(\vartheta)=\frac1{2\pi}
 \sum_{\alpha:\theta_t(\alpha)=\vartheta\;({\rm mod}\;2\pi)}
 \frac{m_t(\alpha)}{|\partial_\alpha\theta_t(\alpha)|}.
 \tag{21.1}
\]

\textbf{Proof.} Around each noncritical preimage, the inverse function
theorem supplies a monotone coordinate chart. Changing variables in
\(m_t(\alpha)d\alpha/(2\pi)\) gives its displayed density contribution.
Summing the charts proves (21.1). On a compact preimage set regularity
makes the number of such charts finite locally. \(\square\)

For a unit probe at angle \(\phi\), the boundaries are \(\phi\pm\pi/2\).
If \(n_t\le B\) throughout the two strips, then

\[
 \Gamma_{\nu_t}(\xi,\eta)\le4B\arcsin\eta.
 \tag{21.2}
\]

If the density is continuous at both boundary angles, then

\[
 \Gamma_{\nu_t}(\xi,\eta)
 =2\arcsin\eta\,[n_t(\phi-\pi/2)+n_t(\phi+\pi/2)]+o(\eta).
 \tag{21.3}
\]

These statements follow by integrating the density over the two arcs,
each of width \(2\arcsin\eta\). A smooth invertible joint
\((\theta,\psi)\) flow alone does not prove any uniform bound on this
projected density. Projection can have folds or large density; a
non-atomicity assertion also needs justification rather than being
inferred from joint invertibility.

There is a weaker but useful connection to angular spread. Restrict the
measure to a fixed measurable label subset \(A\) of mass \(S_A>0\).
Choose a reference angle \(\bar\theta\), define its circular second
moment

\[
 V_A=\frac1{S_A}\int_A d_{\mathbb T}(\theta,\bar\theta)^2d\nu,
\]

and let \(d>0\) be the circular distance from \(\bar\theta\) to the
nearest probe boundary. If \(\arcsin\eta<d\), then

\[
 \boxed{\Gamma_{\nu|A}(\xi,\eta)
 \le\frac{S_AV_A}{(d-\arcsin\eta)^2}.}
 \tag{21.4}
\]

\textbf{Proof.} Every point in the strip is within \(\arcsin\eta\) of a
boundary. The triangle inequality therefore gives
\(d_{\mathbb T}(\theta,\bar\theta)\ge d-\arcsin\eta\). Integrate this
squared inequality over the strip and bound its second moment by the
total second moment. \(\square\)

Equation (21.4) is a finite-tolerance tail bound. For fixed \(V_A\) it
need not vanish as \(\eta\downarrow0\), so it is not a density estimate
and by itself does not prove rate continuity. Nor does input-angle
variance determine the joint input/output covariance in
\(\int ba^\top d\nu\). The same population controls both questions, but
they are not the same observable. A full strip estimate must include the
complement of \(A\), not only a retrospectively named strong family.

For the \(-85^\circ\) probe, the two input gate boundaries are at
\(5^\circ\) and \(185^\circ\). Both must be included. A total
inactive-neuron count is not a measurement of density at either
boundary.

\section{22. Scope guards: finite-horizon persistence, finite width, and
small
initialization}\label{scope-guards-finite-horizon-persistence-finite-width-and-small-initialization}

\subsection{22.1 Rebound does not prove permanent
non-recovery}\label{rebound-does-not-prove-permanent-non-recovery}

After a certified minimum or low-error snapshot \(t_m\), let
\(\Delta_0=E_\xi(t_+)-E_\xi(t_m)>0\). A sufficient finite-horizon
persistence condition is

\[
 \int_{t_+}^{T_{\rm pers}}[-(D_1+D_2)(t,\xi)]_+dt
 \le\Delta_0-\Delta_{\rm pers}
 \quad\text{for some }0<\Delta_{\rm pers}\le\Delta_0.
 \tag{22.1}
\]

Indeed, for every \(t\in[t_+,T_{\rm pers}]\), integrating
\(\dot E=D_1+D_2\) and discarding positive increments gives
\(E_\xi(t)\ge E_\xi(t_m)+\Delta_{\rm pers}\). This is an interface to
prove, not a consequence of weak-family inactivity at an earlier time.

Because \(q>0\), the population Gaussian mixture has strictly positive
density everywhere. Any continuous map with exactly zero population loss
equals the identity everywhere: otherwise a nonzero error at one point
persists on an open neighborhood of positive probability and makes the
loss positive. This observation does not prove convergence to zero loss,
but prevents silently identifying ``very slow recovery'' with ``recovery
is impossible.''

A quantitative conditional version is also available. If the ReLU mass
is at most \(M\) and \(\omega_*=\min_\beta\omega(\beta)>0\), then for a
unit probe

\[
 |e(\xi)|^3\le\frac{8(1+M)\mathcal L}{\omega_*}.
 \tag{22.2}
\]

\textbf{Proof.} Write \(r=|e(\xi)|\le1+M\). On the angular arc of radius
\(r/[2(1+M)]\le1/2\) around the probe, Lipschitz continuity gives
\(|e(a_\beta)|\ge r/2\). The loss is at least
\(\frac12\omega_*[r/(1+M)](r/2)^2\), which yields (22.2). \(\square\)

Thus uniformly bounded mass and \(\mathcal L(t)\to0\) would force
pointwise recovery at every probe. Neither premise is proved here for
long times. The constant \(1/\omega_*\) can be extremely large in
low-density off-cone directions.

The empirical E6 statement in \texttt{ood\_reversals-appendix.tex},
subsection ``Long-horizon persistence and corrected cone geometry,''
concerns ten finite-network seeds through physical time 100. It is
evidence for finite-horizon persistence, not an asymptotic population
theorem.

\subsection{22.2 Exact initial fluctuation scales are not a rate-error
theorem}\label{exact-initial-fluctuation-scales-are-not-a-rate-error-theorem}

For the finite aligned isotropic initialization at fixed total mass
\(S_0\),

\[
 M_h(0)=\frac{S_0}{h}\sum_i a_{\alpha_i}a_{\alpha_i}^\top.
\]

Then

\[
 \mathbb E M_h(0)=\frac{S_0}{2}I_2,\qquad
 \operatorname{sd}([M_h(0)]_{12})=\frac{S_0}{\sqrt{8h}},
\]

\[
 \sqrt{\mathbb E\|M_h(0)-S_0I_2/2\|_{\rm op}^2}
 =\frac{S_0}{2\sqrt h}.
 \tag{22.3}
\]

\textbf{Proof.} Use \(\cos\alpha\sin\alpha=\frac12\sin2\alpha\), whose
mean is zero and second moment is \(1/8\). For the operator norm, write

\[
 M_h-S_0I_2/2=\frac{S_0}{2h}
 \begin{pmatrix}\sum_i\cos2\alpha_i&\sum_i\sin2\alpha_i\\
                 \sum_i\sin2\alpha_i&-\sum_i\cos2\alpha_i\end{pmatrix}.
\]

Its squared operator norm is \(S_0^2/(4h^2)\) times the sum of the two
squared trigonometric sums. Independence gives expectation \(h\).
\(\square\)

The relative root-mean-square operator fluctuation is \(h^{-1/2}\),
about \(7.1\%\) at \(h=200\). It does not imply that the nonlinear
reversal time or signed rates have the same relative error, or that the
continuum equals the finite-seed median.

A sufficient lifting ledger is instead: if
\(|D_p^{(h,N)}-D_p|\le\epsilon_p\) on the relevant windows, individual
population sign margins must exceed their respective \(\epsilon_p\), and
total dominance margins must exceed \(\epsilon_1+\epsilon_2\).
Data-sampling error and probe-gate uncertainty must be included in these
budgets. There is no universal 10-percent cutoff.

\subsection{22.3 Distinguish the two
limits}\label{distinguish-the-two-limits}

At fixed \(S_0\), the continuum limit is \(h\to\infty\) with
\(\varepsilon_h^2=S_0/h\). A fluctuation contribution proportional to
\(\varepsilon_h^2\) then vanishes, whereas an \(S_0\)-scale contribution
need not. A proved lower bound proportional to one of these scales is
not automatically an exact asymptotic law for the full reversal
amplitude.

The separate small-initialization limit is \(S_0\downarrow0\) after
specifying the population model. Its learning times generally move.
Comparing profiles at the same physical time can therefore confound
spread with learning stage. Use a measured, explicitly defined
learning-progress crossing or a stated time shift, and do not assume
monotonicity or uniqueness of that crossing without checking it. Claims
of canonical continuum reversal and of a power law for spread remain
pilot/theorem targets here.

\section{23. Pilot implementation note: eliminate the dense output
matrix}\label{pilot-implementation-note-eliminate-the-dense-output-matrix}

The population pilot is still a numerical witness search. The following
identity reduces its per-evaluation cost without changing the finite
angular quadrature model.

For quadrature labels with reference weights \(w_j^\lambda\), put

\[
 T_j=w_j^\lambda m_j b_{\psi_j}a_{\theta_j}^\top.
\]

For any data angle \(\beta\),

\[
 f(a_\beta)=
 \left[\sum_{j:\theta_j\in(\beta-\pi/2,\beta+\pi/2)\;({\rm mod}\;2\pi)}T_j\right]a_\beta.
 \tag{23.1}
\]

This follows from
\((a_{\theta_j}^\top a_\beta)_+=\mathbf1_{\rm active}a_{\theta_j}^\top a_\beta\).
Sort the current input angles, duplicate them periodically, and form
prefix sums of the four entries of \(T_j\). Each active semicircle is a
prefix-difference query. Thus the output step can cost
\(O(L\log L+K\log L)\), or \(O(L\log L+K+L)\) with a sweep, rather than
\(O(KL)\). This is exact for the chosen discrete label quadrature and
does not require the label-to-input-angle map to remain ordered.

After the data-grid residual is available, evaluate the
moving-half-circle integrals for \(R_p\) with periodic cumulative
integration of \(\omega_p e(a_\beta)a_\beta^\top\). Moving endpoints
need appropriate interpolation or split-cell quadrature. This still has
data-angle discretization error; the prefix identity is not a validation
theorem.

Check data-angle resolution, initial-label resolution, and time step
separately. Record exact rates through both velocity formulas, and log
weighted strip masses at both probe boundaries for a ladder of strip
widths. Also record mass, concept responses, cluster losses, and the
mean/input-output-covariance quantities actually proposed for Theorem B.
Dense direct output evaluation on small random test states remains a
useful implementation cross-check. A scalar exact linear/logistic run
tests normalization independently.

\section{24. Updated proof status}\label{updated-proof-status}

P10 proves all-probe monotonicity and non-harmful cluster contributions
for the matched aligned-isotropic linear population model. P10b proves
immediate creation of a ReLU probe-active off-diagonal response. P11
supplies explicit wrong-gate and angular-tail estimates. P12 gives the
signed Cartesian second-variation formula and a conditional one-sided
tube interface. P13 separates projected strip density from bulk angular
second moments.

These additions do not establish initialized reachability, the canonical
population crossover, substantial weak learning along that trajectory,
or finite-horizon persistence after its rebound. They do not establish
that a coherent closure is sufficient or that covariance is essential.
Those alternatives should be tested with the population pilot and then
resolved by the signed trajectory proof.

\section{25. Additional non-certifying checks for revision
2}\label{additional-non-certifying-checks-for-revision-2}

The new formulas were independently checked with floating-point Gaussian
angular quadrature. These checks did not integrate the initialized
population trajectory and did not run the proposed pilot.

\begin{itemize}
\tightlist
\item
  The explicit \(H(x)\) formula in (17.6) was compared with the original
  half-circle integral at seven vectors in different quadrants; the
  largest absolute discrepancy was \(1.2\cdot10^{-16}\).
\item
  The exact wrong-gate second moment in (18.4) was compared with
  independent one-dimensional Gaussian quadrature at four mean
  orientations; the largest absolute discrepancy was below
  \(9.1\cdot10^{-15}\).
\item
  The gradient-inner-product expression and direct cluster probe
  velocity were compared at a four-label non-balanced state; the largest
  absolute discrepancy was below \(5\cdot10^{-19}\).
\item
  At that same non-balanced state, the three terms of (20.3) were
  approximately \(-0.00231318243\), \(-0.01736117047\), and
  \(-0.00224753606\), summing to \(-0.02192188896\). Centered finite
  differences of the full Cartesian vector field had absolute error
  \(5.8\cdot10^{-10}\) at step \(10^{-3}\), \(5.8\cdot10^{-12}\) at
  \(10^{-4}\), and \(2.7\cdot10^{-13}\) at \(10^{-5}\).
\end{itemize}

The Hessian check intentionally used a non-balanced state, since
Cartesian continuation segments need not preserve balance. These are
transcription and sign/normalization checks, not interval certificates
or substitutes for the proofs.

\end{document}
